\chapter{位元操作}
\label{ch:bitoperations}

位元運算是 C 語言最直接、最高效的操作。 從檢查單一位的狀態，到用位移實
現快速乘除，再到用位元操作交換兩數：這些都是硬體層、密碼學、遊戲引擎常
用的技能。 這章從最基礎的二進位表示開始，逐步講到 5 個核心習語：位元
且、位元或、位元互補、位元互斥或，以及移位。

\section{基礎：進制與二進位表示}
\label{sec:bit-basics}

電腦用二進位儲存資料。 十進位的每一位代表 $10^n$，二進位的每一位代表
$2^n$。 寫下 \texttt{int x = 5;} 時，記憶體裡那一格是
\texttt{00000101}（8 位二進位），也可以寫成 \texttt{0x05}（十六進位）。
後面每一個位元運算子，都是在這一排 $0$ 與 $1$ 上面做事。

\subsection{進制轉換}

C 的整數字面值預設是十進位。 八進位前面加 $0$，十六進位前面加
\texttt{0x}。 印出來時用不同的格式化字串，同一個值可以換成三種進位。

\begin{code}{c}{進制轉換}
int x = 13;      // decimal
// 13 in binary:  1101
// 13 in octal:   015
// 13 in hex:     0x0D

printf("%d\n", x);    // decimal: 13
printf("%o\n", x);    // octal: 15
printf("%x\n", x);    // hex: d
\end{code}

十進位換成二進位，就是把這個數拆成若干個 $2$ 的次方再相加。 $13 = 8 +
4 + 1 = 2^3 + 2^2 + 2^0$，所以二進位是 \texttt{1101}。 反過來，看到
\texttt{1101}，就把亮著的權重加起來：$8 + 4 + 1 = 13$。

後面每一節都要講第 $n$ 位：檢查它、設成 $1$、清成 $0$、切換它。 $n$
從哪一端數、哪一格叫 bit $0$，若不先釘死，\texttt{1u << n} 就對不準
那一格。 下面這張圖把編號、權重、MSB／LSB 一次講完；後文說檢查第
$n$ 位、設置第 $n$ 位，指的都是這個編號。

\subsection{位元的編號、MSB 與權重}

一個位元組有 8 個位，由左到右編號 $7$ 到 $0$。 最左邊是 MSB（most
significant bit，最高有效位），權重最大；最右邊是 LSB（least
significant bit，最低有效位），權重最小。 第 $n$ 位為 $1$ 時，貢獻的
值是 $2^n$。

\begin{center}
\begin{tikzpicture}
    \node[cell] (b7) {0};
    \node[hlcell, right=0pt of b7] (b6) {1};
    \node[cell, right=0pt of b6] (b5) {0};
    \node[cell, right=0pt of b5] (b4) {0};
    \node[hlcell, right=0pt of b4] (b3) {1};
    \node[cell, right=0pt of b3] (b2) {0};
    \node[cell, right=0pt of b2] (b1) {0};
    \node[hlcell, right=0pt of b1] (b0) {1};
    \node[idxlbl, above=2pt of b7] {7};
    \node[idxlbl, above=2pt of b6] {6};
    \node[idxlbl, above=2pt of b5] {5};
    \node[idxlbl, above=2pt of b4] {4};
    \node[idxlbl, above=2pt of b3] {3};
    \node[idxlbl, above=2pt of b2] {2};
    \node[idxlbl, above=2pt of b1] {1};
    \node[idxlbl, above=2pt of b0] {0};
    \node[idxlbl, below=2pt of b7] {$128$};
    \node[idxlbl, below=2pt of b6] {$64$};
    \node[idxlbl, below=2pt of b5] {$32$};
    \node[idxlbl, below=2pt of b4] {$16$};
    \node[idxlbl, below=2pt of b3] {$8$};
    \node[idxlbl, below=2pt of b2] {$4$};
    \node[idxlbl, below=2pt of b1] {$2$};
    \node[idxlbl, below=2pt of b0] {$1$};
\end{tikzpicture}
\diagcaption{上方是位編號，下方是權重 $2^n$。 最左 bit $7$ 是 MSB，最
右 bit $0$ 是 LSB。 橘色格子是值為 $1$ 的位：bit $6$、$3$、$0$，加總
$64 + 8 + 1 = 73$。}
\label{fig:ch12-bit-index}
\end{center}

\texttt{1u << n} 會在第 $n$ 位放一個 $1$，其餘為 $0$，正好對準這一格。
印點陣圖、拆封包時，題目若寫從 MSB 開始，就是從最左邊、編號最大的那
一位先處理。

把 $73$ 拆成權重，過程是：最大不超過 $73$ 的權重是 $64$（bit $6$），
剩下 $9$；再取 $8$（bit $3$），剩下 $1$；最後取 $1$（bit $0$）。 其他
權重都是 $0$。

\keypoint{第 $n$ 位的權重是 $2^n$。 MSB 是最高位，LSB 是最低位；一個
位元組裡 MSB 是 bit $7$，LSB 是 bit $0$。}

\subsection{十六進位：四個位元一組}

二進位寫長了不好對。 十六進位的每一位剛好對應 $4$ 個位元（一個
nibble）：$0$ 是 \texttt{0000}，$9$ 是 \texttt{1001}，$A$ 到 $F$ 是
$10$ 到 $15$。 圖 \ref{fig:ch12-bit-index} 那一列可以切成
\texttt{0100} 與 \texttt{1001}，也就是 $4$ 與 $9$，寫成
\texttt{0x49}。

\begin{example}{c}{(print-73.c) 十進位與十六進位}
#include <stdio.h>

int main(void) {
    unsigned char x = 73;
    printf("dec %u\n", x);
    printf("hex 0x%02x\n", x);
}
\end{example}

\iosample{
}{
dec 73 \\
hex 0x49 \\
}

做位元運算時，遮罩習慣寫十六進位：\texttt{0x0F} 是最低 $4$ 位、
\texttt{0xFF} 是一個位元組、\texttt{0x80} 是 8 位裡的 MSB。 腦子裡再
把每個十六進位數字拆回 $4$ 個位即可。

\subsection{複合賦值運算子}

位元運算常常是拿變數自己跟一個遮罩再運算，再寫回自己。 C 為此準備了簡
寫：\texttt{x op= y} 等價於 \texttt{x = x op y}。 後面的片語會直接用這
種寫法。

\begin{code}{c}{複合賦值}
x &= mask;    // same as x = x & mask
x |= bit;     // same as x = x | bit
x ^= bit;     // same as x = x ^ bit
x <<= n;      // same as x = x << n
x >>= n;      // same as x = x >> n
\end{code}

五種位元運算子都有對應的複合賦值。 寫成 \texttt{x |= (1u << 2)} 的時
候，心裡仍然要把它展開成：先算出右邊，再跟左邊做 OR，結果放回
\texttt{x}。

\subsection{有符號數的二補數}

最高位（MSB）在有號整數裡同時是符號位：$0$ 表示非負，$1$ 表示負。 現
代 C 用二補數（two's complement）表示負數。

\begin{example}{c}{(twos-complement.c) 補碼表示}
#include <stdio.h>

int main(void) {
    signed char x = 5;     // binary: 00000101
    signed char neg_x = -x; // binary: 11111011

    printf("5 in hex: 0x%02x\n", (unsigned char)x);
    printf("-5 in hex: 0x%02x\n", (unsigned char)neg_x);
}
\end{example}

\iosample{
}{
5 in hex: 0x05 \\
-5 in hex: 0xfb \\
}

取負的規則可以拆成兩步：先把每一位翻反（\texttt{\~{}x}），再加 $1$。
\texttt{5} 是 \texttt{00000101}，翻反得 \texttt{11111010}，加 $1$ 得
\texttt{11111011}，也就是範例印出的 \texttt{0xfb}。 後面講位元互補時，
會再用到這個翻反的動作。

二補數讓加法本身就能處理正負：$5 + (-5)$ 在 8 位裡是
\texttt{00000101} $+$ \texttt{11111011}，加上進位後低 8 位是
\texttt{00000000}。 $-1$ 的每一位都是 $1$（\texttt{0xff}）：把
\texttt{00000000} 翻反再加 $1$，進位走出格外，留下 8 個 $1$。

\begin{example}{c}{(minus-one.c) -1 與相加歸零}
#include <stdio.h>

int main(void) {
    signed char n = -1;
    signed char a = 5;
    signed char b = -5;
    printf("-1 in hex: 0x%02x\n", (unsigned char)n);
    printf("5 + (-5) = %d\n", a + b);
}
\end{example}

\iosample{
}{
-1 in hex: 0xff \\
5 + (-5) = 0 \\
}

8 位有號的範圍是 $-128$ 到 $127$：\texttt{10000000} 是 $-128$，沒有對
應的正數可再取負。 無號的 \texttt{\~{}0u} 與有號的 $-1$ 位元模式相同
，混在一起時正負會看錯。

\keypoint{二補數取負就是 \texttt{\~{}x + 1}。 $-1$ 的每一位都是 $1$；
有號整數的 MSB 同時是符號位。}

\stylepoint{做位元運算時，變數請宣告成 \texttt{unsigned}。 有號數的右
移、把 $1$ 左移進符號位，行為要嘛未定義、要嘛依實作而定；
\texttt{unsigned} 讓每一位都只是 $0$ 或 $1$，沒有符號位延伸這條後路。}

\subsection{五個運算子怎麼算、怎麼用}

位元運算是對每一個位各自計算，左右鄰居互不干擾，也沒有進位。 這點跟
二補數加法不同：加法會把進位傳到下一格，AND／OR／XOR 不會。

\begin{center}
\begin{tabular}{|l|l|l|}
\hline
\textbf{運算子} & \textbf{每一位} & \textbf{典型用途} \\
\hline
\texttt{\&} & 兩位都是 $1$ 才得 $1$ & 檢查、遮罩 \\
\texttt{|} & 有一位是 $1$ 就得 $1$ & 設置成 $1$ \\
\texttt{\~{}} & $0$ 變 $1$、$1$ 變 $0$ & 做清除用的擋板 \\
\texttt{\textasciicircum{}} & 不同得 $1$、相同得 $0$ & 切換、交換 \\
\texttt{<<}／\texttt{>>} & 整排左右滑 & 對準某一格、乘除 $2^n$ \\
\hline
\end{tabular}
\end{center}

手算時把兩個運算元按 bit $7$ 到 bit $0$ 上下對齊，逐欄套真值表。 後面
五節各對準表裡的一行：且用來檢查、或用來設置、互補用來清除、互斥或用
來交換、移位用來對準再測試或計數。 遮罩這塊擋板會一路跟著走：AND 用
它留下洞、OR 往洞裡塞 $1$、NOT 把洞翻成擋板再清位、XOR 只翻洞裡那一
格。 章末再用同一組旗標與一組檔案權限，把五個運算子串成一支小程式。

\begin{drills}[label=drills:bit-basics]
\item 把十進位 $13$ 拆成 $2$ 的次方相加，寫出對應的二進位。
\item 一個位元組的 MSB 是第幾位？ 它的權重是多少？
\item \texttt{0x49} 換成二進位是多少？ 換成十進位呢？
\item \texttt{x |= (1u << 2)} 這種複合賦值寫法，展開後等價於哪一行完
整算式？
\end{drills}

\section{位元且：檢查位元}
\label{sec:bit-and}

\texttt{x \& y} 對每一位做邏輯 AND。 最常用的是檢查特定位是否為 1。

\keypoint{位元且的真值表：0 AND 0 $= 0$，0 AND 1 $= 0$，1 AND 0 $= 0$，
1 AND 1 $= 1$。 只有兩位都是 1 才得 1。}

\subsection[逐位元計算：AND]{逐位元計算}

AND 是對每一個位各自運算，左右鄰居互不干擾。 把兩個運算元上下對齊，逐
欄看過一次，就能用手算出結果。 橘色格子是這一列為 $1$ 的位；結果列只
在兩排都是橘色的欄留下 $1$。

\begin{center}
\begin{tikzpicture}
    \foreach \col/\bn in {0/7,1/6,2/5,3/4,4/3,5/2,6/1,7/0} {
        \node[idxlbl] at (0.75*\col, 1.35) {\bn};
    }
    \node[idxlbl, anchor=east] at (-0.45, 0.75) {147};
    \foreach \col/\v in {0/1,1/0,2/0,3/1,4/0,5/0,6/1,7/1} {
        \ifnum\v=1
            \node[hlcell, minimum width=0.75cm, minimum height=0.75cm]
                at (0.75*\col, 0.75) {\v};
        \else
            \node[cell, minimum width=0.75cm, minimum height=0.75cm]
                at (0.75*\col, 0.75) {\v};
        \fi
    }
    \node[idxlbl, anchor=east] at (-0.45, 0) {\& 61};
    \foreach \col/\v in {0/0,1/0,2/1,3/1,4/1,5/1,6/0,7/1} {
        \ifnum\v=1
            \node[hlcell, minimum width=0.75cm, minimum height=0.75cm]
                at (0.75*\col, 0) {\v};
        \else
            \node[cell, minimum width=0.75cm, minimum height=0.75cm]
                at (0.75*\col, 0) {\v};
        \fi
    }
    \node[idxlbl, anchor=east] at (-0.45, -0.75) {17};
    \foreach \col/\v in {0/0,1/0,2/0,3/1,4/0,5/0,6/0,7/1} {
        \ifnum\v=1
            \node[hlcell, minimum width=0.75cm, minimum height=0.75cm]
                at (0.75*\col, -0.75) {\v};
        \else
            \node[cell, minimum width=0.75cm, minimum height=0.75cm]
                at (0.75*\col, -0.75) {\v};
        \fi
    }
\end{tikzpicture}
\diagcaption{$147 \& 61$。 上下對齊後逐欄做 AND，只有 bit $4$ 與
bit $0$ 兩排都是 $1$，結果是 $16 + 1 = 17$。}
\label{fig:ch12-and-align}
\end{center}

\begin{code}{c}{逐位元手算 AND}
// bit:  7 6 5 4 3 2 1 0
//       1 0 0 1 0 0 1 1   (147)
//     & 0 0 1 1 1 1 0 1   (61)
//       ---------------
//       0 0 0 1 0 0 0 1   (17)  // bits 4 and 0 survive
\end{code}

十進位驗算：$147 \& 61 = 17$。 後面每次看到 \texttt{\&}，都可以在腦子
裡排成這樣的兩行再對一次。 遮罩就是下面那一排：為 $1$ 的欄是擋板上的
洞，為 $0$ 的欄把上面那一排蓋掉。

\subsection{檢查第 n 位}

要檢查第 $n$ 位是否為 $1$，先把 \texttt{1u} 左移 $n$ 格做出遮罩，再跟
原數做 AND。

\begin{code}{c}{檢查第 n 位}
unsigned int x = 12;  // binary: 1100
int bit2 = (x >> 2) & 1;  // shift right 2, check lowest bit
printf("%d\n", bit2);  // output: 1

// or: AND with a 1u mask
int bit3 = (x & (1u << 3)) != 0;  // check bit 3
printf("%d\n", bit3);  // output: 1
\end{code}

最低位（bit $0$）為 $1$ 就是奇數。 \texttt{x \& 1} 只開這一格的洞，結
果不是 $0$ 就是 $1$。

\begin{code}{c}{用最低位檢查奇偶}
int x = 7;
if (x & 1) {
    printf("odd\n");   // 7 AND 1 = 1
} else {
    printf("even\n");
}
\end{code}

\subsection{應用：檢查多個位}

遮罩（mask）是預先設置好的位元組合，用來篩選特定位。

\begin{code}{c}{多位元遮罩}
unsigned int status = 0b1101;  // 13
unsigned int mask = 0b1010;    // 10

int result = status & mask;  // AND: only bit 3 and bit 1 pass
printf("0x%x\n", result);  // output: 0x8
\end{code}

\subsection{遮罩的意義}

遮罩可以想成一塊擋板，上面打了幾個洞。 洞的位置（位元是 $1$）讓原來的
位看得見、原樣通過；擋板擋住的位置（位元是 $0$）強迫結果變成 $0$。 所
以 \texttt{x \& mask} 只留下 mask 為 $1$ 的那些位，其餘一律清掉。

同一個想法也可以用來截斷。 \texttt{ch \&= 0xFF;} 只開最低 8 個洞，高
於一個位元組的位全部被擋住，結果保證落在 $0$ 到 $255$。 從較寬的整數
裡取出最低一個 byte 時，這行很常見。

\subsection[陷阱：flags == MASK 比較法錯了]{陷阱：不要直接比較 \texttt{flags == MASK}}

要確認 MASK 裡的每一位都亮著，條件寫成
\texttt{(flags \& MASK) == MASK}。 這個寫法只看 MASK 覆蓋到的那些位；
\texttt{flags} 上其他位是 $0$ 或 $1$ 都不影響結果。

\begin{code}{c}{用遮罩檢查旗標}
unsigned int flags = 0x7;   // bits 0, 1, 2
unsigned int MASK  = 0x5;   // bits 0 and 2

if (flags == MASK) {              // 7 == 5, false
    printf("exact match\n");
}
if ((flags & MASK) == MASK) {     // (7 & 5) == 5, true
    printf("mask bits on\n");
}
\end{code}

\texttt{flags} 是 $7$，\texttt{MASK} 是 $5$。 直接比較兩個整數，會要求
每一位都相同，多開的 bit $1$ 讓比較失敗。 先做 AND 再比，只問 bit $0$
與 bit $2$ 有沒有亮，答案是有。

括號一定要寫。 \texttt{\&} 的優先序低於 \texttt{==}，若寫成
\texttt{flags \& MASK == MASK}，C 會先算 \texttt{MASK == MASK}（結果是
$1$），再跟 \texttt{flags} 做 AND，整個條件的意義就變了。

\subsection{具名旗標：檢查}

下一節會用 \texttt{FLAG\_BLUE}、\texttt{FLAG\_GREEN}、
\texttt{FLAG\_RED} 來設置旗標。 檢查走的就是這一節的 AND：洞開在要問
的那幾位，其餘擋住。

\begin{example}{c}{(flag-check.c) 檢查具名旗標}
#include <stdio.h>

#define FLAG_BLUE  1u
#define FLAG_GREEN 2u
#define FLAG_RED   4u

int main(void) {
    unsigned int flags = FLAG_BLUE | FLAG_RED;

    if ((flags & FLAG_BLUE) != 0) {
        printf("blue on\n");
    }
    if ((flags & FLAG_GREEN) != 0) {
        printf("green on\n");
    } else {
        printf("green off\n");
    }
    if ((flags & (FLAG_BLUE | FLAG_RED)) == (FLAG_BLUE | FLAG_RED)) {
        printf("blue and red on\n");
    }
}
\end{example}

\iosample{
}{
blue on \\
green off \\
blue and red on \\
}

\texttt{FLAG\_BLUE | FLAG\_RED} 先合成要問的那兩格。 跟
\texttt{flags == MASK} 那個陷阱同一句：先 AND 再比 MASK，多開的位不影
響。

\subsection{練習：手算 AND}

把 $182$（\texttt{10110110}）跟 $205$（\texttt{11001101}）上下對齊，
逐欄做 AND。 寫完再對解答：$10000100$，也就是 $132$。 只有 bit $7$
與 bit $2$ 兩排都是 $1$。

AND 留下擋板上的洞。 下一節的 OR 則是往指定的格子裡塞 $1$，已經是
$1$ 的位維持 $1$。

\begin{drills}[label=drills:bit-and]
\item 手算 $147 \& 61$（bit 7 到 bit 0 上下對齊逐欄算），結果是多少？
\item 要測試變數 \texttt{x} 的第 3 位是否為 1，該怎麼寫？
\item 為什麼 \texttt{(flags \& MASK) == MASK} 是對的，而
\texttt{flags == MASK} 通常是錯的？
\item \texttt{flags \& MASK == MASK}（沒加括號）會被剖析成什麼？ 為什
麼跟原意不同？
\end{drills}

\section{位元或：設置位元}
\label{sec:bit-or}

\texttt{x | y} 對每一位做邏輯 OR。 最常用的是把特定位設為 1。 手算時
一樣上下對齊：任一欄有 $1$，結果就是 $1$。 跟 AND 那張圖比，這裡只要
有一排是橘色，結果列就是橘色。

\begin{center}
\begin{tikzpicture}
    \foreach \col/\bn in {0/7,1/6,2/5,3/4,4/3,5/2,6/1,7/0} {
        \node[idxlbl] at (0.75*\col, 1.35) {\bn};
    }
    \node[idxlbl, anchor=east] at (-0.45, 0.75) {128};
    \foreach \col/\v in {0/1,1/0,2/0,3/0,4/0,5/0,6/0,7/0} {
        \ifnum\v=1
            \node[hlcell, minimum width=0.75cm, minimum height=0.75cm]
                at (0.75*\col, 0.75) {\v};
        \else
            \node[cell, minimum width=0.75cm, minimum height=0.75cm]
                at (0.75*\col, 0.75) {\v};
        \fi
    }
    \node[idxlbl, anchor=east] at (-0.45, 0) {| 4};
    \foreach \col/\v in {0/0,1/0,2/0,3/0,4/0,5/1,6/0,7/0} {
        \ifnum\v=1
            \node[hlcell, minimum width=0.75cm, minimum height=0.75cm]
                at (0.75*\col, 0) {\v};
        \else
            \node[cell, minimum width=0.75cm, minimum height=0.75cm]
                at (0.75*\col, 0) {\v};
        \fi
    }
    \node[idxlbl, anchor=east] at (-0.45, -0.75) {132};
    \foreach \col/\v in {0/1,1/0,2/0,3/0,4/0,5/1,6/0,7/0} {
        \ifnum\v=1
            \node[hlcell, minimum width=0.75cm, minimum height=0.75cm]
                at (0.75*\col, -0.75) {\v};
        \else
            \node[cell, minimum width=0.75cm, minimum height=0.75cm]
                at (0.75*\col, -0.75) {\v};
        \fi
    }
\end{tikzpicture}
\diagcaption{$128 | 4$。 bit $7$ 與 bit $2$ 各有一排是 $1$，結果兩格
都亮，得到 $132$。}
\label{fig:ch12-or-align}
\end{center}

\begin{code}{c}{逐位元手算 OR}
// bit:  7 6 5 4 3 2 1 0
//       1 0 0 0 0 0 0 0   (128)
//     | 0 0 0 0 0 1 0 0   (4)
//       ---------------
//       1 0 0 0 0 1 0 0   (132)
\end{code}

\subsection{設置第 n 位為 1}

遮罩寫成 \texttt{1u << n}。 字面值 \texttt{1} 的型別是有號
\texttt{int}；在 $32$ 位系統上，\texttt{1 << 31} 會把 $1$ 推進符號位，
行為沒有定義。 \texttt{1u} 是 \texttt{unsigned}，\texttt{1u << 31} 得
到 \texttt{0x80000000u}，每一位的意義都穩定。

\begin{code}{c}{設置第 n 位}
unsigned int x = 8;   // binary: 1000
x |= (1u << 2);      // set bit 2: OR with 0100
printf("0x%x\n", x);  // output: 0xc (12 in decimal)
\end{code}

\subsection{應用：開啟功能}

在嵌入式系統中，硬體暫存器用位表示開關狀態。

\begin{code}{c}{硬體暫存器}
unsigned int flags = 0;
flags |= (1u << 0);  // enable interrupt
flags |= (1u << 2);  // enable timer
printf("0x%x\n", flags);  // output: 0x5
\end{code}

\subsection{具名旗標}

把每個位的意義寫成巨集，讀程式的人就不必心算第幾位代表什麼。 習慣上每
個旗標只占一個位，值是 $1$、$2$、$4$、$8$ 這種 $2$ 的次方，互相不會重疊。

\begin{code}{c}{具名旗標}
#define FLAG_BLUE  1    // bit 0
#define FLAG_GREEN 2    // bit 1
#define FLAG_RED   4    // bit 2

unsigned int flags = 0;
flags |= FLAG_BLUE | FLAG_RED;   // turn on blue and red at once
// flags is now 5 (101)
\end{code}

\texttt{FLAG\_BLUE | FLAG\_RED} 先合成遮罩 $5$（bit $0$ 與 bit $2$），
再 OR 進 \texttt{flags}。 已經亮著的位維持亮著，原本是 $0$ 的位被設成
$1$。 上一節用同一組名字做檢查，下一節用它們做清除。

\subsection{陷阱：\texttt{=} 會蓋掉其他旗標}

要把藍色打開，正確是 \texttt{flags |= FLAG\_BLUE}。 若寫成
\texttt{flags = FLAG\_BLUE}，整格被換成 $1$，原本亮著的綠、紅全部熄
掉。 OR 只把遮罩為 $1$ 的位點亮，其餘位原樣留下；指派則是整份覆寫。

\begin{code}{c}{設置用 OR，不要整份覆寫}
unsigned int flags = FLAG_GREEN | FLAG_RED;  // 6
flags = FLAG_BLUE;     // now 1; green and red are gone
flags = FLAG_GREEN | FLAG_RED;
flags |= FLAG_BLUE;    // now 7; green and red stay
\end{code}

\subsection{練習：手算 OR}

同一對數字：$182$（\texttt{10110110}）OR $205$（\texttt{11001101}）。
解答是 $11111111$，也就是 $255$。 每一欄至少有一個 $1$。

OR 把指定的位設成 $1$。 要關回去，得先做一塊只擋那幾格的板，再 AND
回去，就是下一節的互補。

\keypoint{位元或的真值表：任何位 OR 0 不變，任何位 OR 1 變 1。}

\begin{drills}[label=drills:bit-or]
\item 手算 $182 | 205$（二進位對齊逐欄算），結果是多少？
\item 為什麼設置某個旗標要用 \texttt{|=}，直接寫 \texttt{flags =
FLAG\_BLUE} 會有什麼後果？
\item 為什麼遮罩要寫成 \texttt{1u << n} 而不是 \texttt{1 << n}？
\end{drills}

\section{位元互補：清除位元}
\label{sec:bit-complement}

\texttt{\~{} x} 翻反 \texttt{x} 的每一位：0 變 1，1 變 0。 最常用的是把
特定位清為 0。

\subsection{清除第 n 位為 0}

\begin{code}{c}{清除第 n 位}
unsigned int x = 12;   // binary: 1100
x &= ~(1u << 2);      // clear bit 2: AND with ~0100 = 1011
printf("0x%x\n", x);  // output: 0x8 (8 in decimal)
\end{code}

注意：\texttt{\~{}(1u << 2)} 先左移再翻反。 括號很重要。

\subsection{應用：關閉功能}

\begin{code}{c}{關閉功能}
unsigned int flags = 0x7;  // binary: 0111
flags &= ~(1u << 1);      // disable bit 1
printf("0x%x\n", flags);  // output: 0x5 (0101)
\end{code}

\subsection{清除具名旗標}

延續上一節的 \texttt{FLAG\_*}。 要同時關掉好幾個旗標，先用 OR 把它們
合成一塊，再取互補，最後 AND 回去。

\begin{code}{c}{清除具名旗標}
#define FLAG_BLUE  1
#define FLAG_GREEN 2
#define FLAG_RED   4

unsigned int flags = FLAG_BLUE | FLAG_GREEN | FLAG_RED;  // 7
flags &= ~(FLAG_GREEN | FLAG_RED);   // turn off green and red
// flags is now FLAG_BLUE (1)
\end{code}

先算 \texttt{FLAG\_GREEN | FLAG\_RED}，得到 $6$（二進位 \texttt{110}）。
\texttt{\~{}6} 把這兩位翻成 $0$、其餘翻成 $1$。 再跟原本的 $7$
（\texttt{111}）做 AND，只留下 bit $0$，也就是 \texttt{FLAG\_BLUE}。
補碼那一節說過取負是 \texttt{\~{}x + 1}；這裡只用到前半段的翻反，用來
做擋板。

\subsection[可攜寫法：全部是 1 用 ~0]{可攜寫法：\texttt{\~{}0} 代表全部是 1}

把遮罩寫成 \texttt{0xFF} 或 \texttt{0xFFFF}，等於把位元寬度綁死在 $8$
位或 $16$ 位。 \texttt{int} 若是 $32$ 位，\texttt{0xFF} 只開最低 8 個
洞，高位全是 $0$。 換成 \texttt{\~{}0}，每一位都會是 $1$，寬度跟著當
前型別走。

unsigned 語境請寫 \texttt{\~{}0u}。 有號的 \texttt{\~{}0} 在補碼裡就是
$-1$，跟整數運算混在一起時容易看錯正負。

\begin{code}{c}{可攜的全 1 遮罩}
unsigned int x = 0x3F;     // 00111111
x &= ~0x1Fu;               // clear the lowest 5 bits
// x is now 0x20 (00100000)
\end{code}

\texttt{0x1F} 是最低 5 位為 $1$。 \texttt{\~{}0x1Fu} 把這 5 位翻成 $0$
、其餘維持 $1$，再 AND 回去就清掉最低 5 位。
\texttt{0x3F \& \~{}0x1F} 得 \texttt{0x20}，跟手算一致。 這裡沒有寫死
\texttt{0xFFFFFFE0} 這種跟 $32$ 位綁在一起的常數。

清除是 AND 加上一塊翻反過的擋板：洞變成牆，牆變成洞。 若只要翻某一格
、其他位不動，擋板的想法仍然適用，運算子改成 XOR，見下一節。

\keypoint{位元互補的真值表：NOT 0=1，NOT 1=0。 翻反所有位。
常與 AND 搭配使用來清除特定位。}

\begin{drills}[label=drills:bit-complement]
\item \texttt{x = 12}（\texttt{1100}），執行 \texttt{x \&= \~{}(1u <<
2)} 之後 \texttt{x} 是多少？
\item 為什麼要寫 \texttt{\~{}0u} 而不是寫死 \texttt{0xFF} 或
\texttt{0xFFFF} 當全 1 遮罩？
\item 要同時清除 \texttt{FLAG\_GREEN} 與 \texttt{FLAG\_RED} 兩個旗標，
該怎麼寫成一行？
\end{drills}

\section{位元互斥或：切換與交換}
\label{sec:bit-xor}

\texttt{x \textasciicircum{} y} 對每一位做邏輯 XOR。 相同為 $0$，不同
為 $1$。 最常用的是切換某一格，以及不用臨時變數交換兩個整數。 這兩件
事都靠同一條性質：同一個值 XOR 兩次，會回到原值。

\subsection[逐位元計算：XOR]{逐位元計算}

XOR 一樣上下對齊、逐欄算。 兩排相同的欄得 $0$，不同的欄得 $1$。

\begin{center}
\begin{tikzpicture}
    \foreach \col/\bn in {0/7,1/6,2/5,3/4,4/3,5/2,6/1,7/0} {
        \node[idxlbl] at (0.75*\col, 1.35) {\bn};
    }
    \node[idxlbl, anchor=east] at (-0.45, 0.75) {5};
    \foreach \col/\v in {0/0,1/0,2/0,3/0,4/0,5/1,6/0,7/1} {
        \ifnum\v=1
            \node[hlcell, minimum width=0.75cm, minimum height=0.75cm]
                at (0.75*\col, 0.75) {\v};
        \else
            \node[cell, minimum width=0.75cm, minimum height=0.75cm]
                at (0.75*\col, 0.75) {\v};
        \fi
    }
    \node[idxlbl, anchor=east] at (-0.45, 0) {\textasciicircum{} 12};
    \foreach \col/\v in {0/0,1/0,2/0,3/0,4/1,5/1,6/0,7/0} {
        \ifnum\v=1
            \node[hlcell, minimum width=0.75cm, minimum height=0.75cm]
                at (0.75*\col, 0) {\v};
        \else
            \node[cell, minimum width=0.75cm, minimum height=0.75cm]
                at (0.75*\col, 0) {\v};
        \fi
    }
    \node[idxlbl, anchor=east] at (-0.45, -0.75) {9};
    \foreach \col/\v in {0/0,1/0,2/0,3/0,4/1,5/0,6/0,7/1} {
        \ifnum\v=1
            \node[hlcell, minimum width=0.75cm, minimum height=0.75cm]
                at (0.75*\col, -0.75) {\v};
        \else
            \node[cell, minimum width=0.75cm, minimum height=0.75cm]
                at (0.75*\col, -0.75) {\v};
        \fi
    }
\end{tikzpicture}
\diagcaption{$5 \oplus 12 = 9$。 bit $3$ 與 bit $0$ 兩排不同，結果這
兩格是 $1$；其餘欄相同，結果是 $0$。}
\label{fig:ch12-xor-align}
\end{center}

\keypoint{位元互斥或的真值表：0 XOR 0=0，0 XOR 1=1，1 XOR 0=1，1 XOR 1=0。
相同變 0，不同變 1。 任何數 XOR 自己都是 $0$；任何數 XOR $0$ 都不變。}

\subsection{同一值 XOR 兩次會抵銷}

對每一位來說，把同一個 $k$ 的位 XOR 進去兩次，等於沒做。 寫成式子：
\texttt{(x \textasciicircum{} k) \textasciicircum{} k} 等於
\texttt{x}。 先記住這條，下面的交換與加密都是它的應用。

\begin{code}{c}{XOR 兩次回到原值}
unsigned int x = 0x48;   // 'H'
unsigned int k = 0x2A;
unsigned int y = x ^ k;
unsigned int back = y ^ k;   // back == x
\end{code}

\subsection{為什麼 XOR 交換行得通}

令一開始 $a$ 存 $a_0$、$b$ 存 $b_0$。 三步之後，$a$ 應變成 $b_0$，
$b$ 應變成 $a_0$。 每一步都用上一小節的抵銷。

\begin{center}
\begin{tabular}{|l|l|l|l|}
\hline
\textbf{步驟} & \textbf{$a$} & \textbf{$b$} & \textbf{這一步做了什麼} \\
\hline
開始 & $a_0$（$5$，\texttt{0101}） & $b_0$（$12$，\texttt{1100}） &
兩格各自的原值 \\
\texttt{a \textasciicircum{}= b} & $a_0 \oplus b_0$（$9$，\texttt{1001}） &
$b_0$ & $a$ 暫時變成兩者的 XOR \\
\texttt{b \textasciicircum{}= a} & $a_0 \oplus b_0$ &
$b_0 \oplus (a_0 \oplus b_0) = a_0$ & $b_0$ 被 XOR 兩次，抵銷，留下
$a_0$ \\
\texttt{a \textasciicircum{}= b} &
$(a_0 \oplus b_0) \oplus a_0 = b_0$ & $a_0$ & $a_0$ 被 XOR 兩次，抵
銷，留下 $b_0$ \\
\hline
\end{tabular}
\end{center}

第三步的 $b$ 已經是 $a_0$，所以最後一次 XOR 等於把 $a$ 裡的 $a_0$ 消
掉，只剩 $b_0$。

\begin{code}{c}{用 XOR 交換}
int a = 5;   // binary: 0101
int b = 12;  // binary: 1100

a = a ^ b;  // a: 0101 XOR 1100 = 1001 (9)
b = b ^ a;  // b: 1100 XOR 1001 = 0101 (5)
a = a ^ b;  // a: 1001 XOR 0101 = 1100 (12)

printf("a=%d, b=%d\n", a, b);  // output: a=12, b=5
\end{code}

這個技巧只對整數有效，不會用到額外記憶體。

\subsection[陷阱：交換的兩變數不能同位址]{陷阱：交換的兩個變數不能是同一個位址}

XOR 交換法要求兩個運算元是不同的變數。 任何數 XOR 自己都是 $0$；若對
同一個變數連續做三次 \texttt{a = a \textasciicircum{} a}，第一步就已經
把資料清成 $0$，後面兩步只是 $0$ XOR $0$。

\begin{code}{c}{對同一變數做 XOR 交換}
int a = 5;
a = a ^ a;   // 5 XOR 5 = 0
a = a ^ a;   // 0 XOR 0 = 0
a = a ^ a;   // still 0
printf("%d\n", a);   // output: 0
\end{code}

呼叫函式時若兩個參數其實指向同一格記憶體，例如 \texttt{swap(\&x, \&x)}
，也會踩進這個陷阱。 日常交換兩個整數，用一個臨時變數比較不容易寫錯。

\subsection{具名旗標：切換}

OR 把位設成 $1$，AND 加互補把位清成 $0$。 XOR 則是亮的變暗、暗的變亮
，做一次切一次，做兩次回到原狀，正好是抵銷性質。

\begin{example}{c}{(flag-toggle.c) 切換具名旗標}
#include <stdio.h>

#define FLAG_BLUE  1u
#define FLAG_GREEN 2u
#define FLAG_RED   4u

int main(void) {
    unsigned int flags = FLAG_BLUE | FLAG_RED;
    printf("start: %u\n", flags);
    flags ^= FLAG_GREEN;
    printf("toggle green on: %u\n", flags);
    flags ^= FLAG_GREEN;
    printf("toggle green off: %u\n", flags);
}
\end{example}

\iosample{
}{
start: 5 \\
toggle green on: 7 \\
toggle green off: 5 \\
}

\texttt{FLAG\_GREEN} 只占 bit $1$。 第一次 XOR 把它點亮，$5$ 變成
$7$；第二次再 XOR 同一格，綠燈熄掉，回到 $5$。 藍與紅沒被碰到。

\subsection{練習：手算 XOR}

同一對數字：$182$（\texttt{10110110}）XOR $205$（\texttt{11001101}）。
解答是 $01111011$，也就是 $123$。 跟前面 AND、OR 的練習比：AND 只留兩
排都是 $1$ 的欄，OR 留至少一排是 $1$ 的欄，XOR 留剛好一排是 $1$ 的欄。

切換與交換都靠 XOR 的可逆性。 要對準某一格、或把幾個位元組拆開再組回
去，還需要把整排左右滑，見下一節。

\subsection{應用：XOR 加密}

交換法用過的同一條性質也可以拿來加密。 明文跟一把固定的 key 做 XOR，
得到密文；密文再跟同一把 key 做 XOR，明文就回來了。

\begin{example}{c}{(xor-cipher.c) XOR 加密}
#include <stdio.h>

int main(void) {
    unsigned char msg0 = 'H';
    unsigned char msg1 = 'I';
    unsigned char key = 0x2A;

    unsigned char enc0 = msg0 ^ key;
    unsigned char enc1 = msg1 ^ key;
    unsigned char dec0 = enc0 ^ key;
    unsigned char dec1 = enc1 ^ key;

    printf("plain: %c%c\n", msg0, msg1);
    printf("cipher: 0x%02x 0x%02x\n", enc0, enc1);
    printf("decoded: %c%c\n", dec0, dec1);
}
\end{example}

\iosample{
}{
plain: HI \\
cipher: 0x62 0x63 \\
decoded: HI \\
}

\texttt{'H'} 是 \texttt{0x48}，\texttt{'I'} 是 \texttt{0x49}，key 是
\texttt{0x2A}。 \texttt{0x48 \textasciicircum{} 0x2A} 得 \texttt{0x62}，
\texttt{0x49 \textasciicircum{} 0x2A} 得 \texttt{0x63}。 密文再 XOR 一
次 \texttt{0x2A}，兩個位元組都回到明文。 這把固定 key 的作法只適合用
來理解 XOR 的可逆性；真正的加密還需要更長、不重複使用的金鑰。

\begin{drills}[label=drills:bit-xor]
\item 手算 $182 \oplus 205$，結果是多少？
\item 為什麼 \texttt{(x \textasciicircum{} k) \textasciicircum{} k}
會等於 \texttt{x}？
\item XOR 交換法為什麼不能用在 \texttt{swap(\&x, \&x)} 這種情況？
\item \texttt{flags \textasciicircum{}= FLAG\_GREEN} 連續執行兩次，
\texttt{flags} 最後會變成什麼？
\end{drills}

\section{移位：數字移位、測試與計數}
\label{sec:bit-shift}

前面四節都假設遮罩已經對準要動的那一格。 對準的動作就是移位：
\texttt{1u << n} 做出第 $n$ 位的擋板，\texttt{x >> n} 把第 $n$ 位滑到
最低位再檢查。 這一節把移位本身講清楚，並用它拆、組多位元組的整數。

\subsection{左移與右移}

\texttt{x << n} 左移 n 格，等同乘以 $2^n$。 \texttt{x >> n} 右移 n 格，
等同除以 $2^n$（整數除法）。

\begin{code}{c}{左移與右移}
unsigned int x = 3;   // binary: 11
unsigned int y = x << 2; // left shift: 1100 = 12
unsigned int z = y >> 1; // right shift: 0110 = 6

printf("%u << 2 = %u\n", x, y);  // 3 << 2 = 12
printf("%u >> 1 = %u\n", y, z);  // 12 >> 1 = 6
\end{code}

\keypoint{十進位裡把小數點右移兩位，等於乘以 $100$。 二進位的左移 $n$
位，同樣是乘以 $2^n$。 這個類比在非負數與 \texttt{unsigned} 上完全成
立。}

\subsection{有號數右移：算術位移與邏輯位移}

右移空出來的高位要補什麼，得看型別。 \texttt{unsigned} 永遠補 $0$，這
叫邏輯位移。 有號負數的右移在 C 裡是 implementation-defined：編譯器可
以補 $0$，也可以複製符號位（算術位移，負數會一直保持負的）。

同一份程式換一台機器、換一個編譯器，有號右移的結果可能跟著變。 可攜的
寫法是：有號數若只是要除以 $2$ 的次方，直接用 \texttt{/}，不要用
\texttt{>>}。

就算編譯器對負數做算術位移，結果也跟 \texttt{/} 不同。
\texttt{-3 / 2} 得到 $-1$（向零取整）；算術右移一位通常得到 $-2$（向
負無窮取整）。 兩條路連取整的方向都不一樣，混用會悄悄算錯。

\subsection{測試位元（用右移）}

最直接的方法：右移目標位到最低位，再用 AND 檢查。

\begin{code}{c}{用右移測試位元}
unsigned int x = 12;  // binary: 1100
int bit2 = (x >> 2) & 1;  // shift right 2, check bit
printf("bit 2: %d\n", bit2);  // output: 1

int bit0 = (x >> 0) & 1;  // check bit 0
printf("bit 0: %d\n", bit0);  // output: 0
\end{code}

\subsection{應用：從整數中取出色彩分量}

一個 \texttt{unsigned long} 可以一次塞進三個 byte。 常見的包裝是把紅、
綠、藍依序放進不同的 8 位區間：最高那個 byte 是 R，中間是 G，最低是 B。
要拆開來用，先右移對準，再跟 \texttt{0xFF} 做 AND，只留下一個位元組。

\begin{example}{c}{(extract-color.c) 取出 RGB 分量}
#include <stdio.h>

int main(void) {
    unsigned long color = 0xC04080UL;
    unsigned int r = (color >> 16) & 0xFF;
    unsigned int g = (color >> 8) & 0xFF;
    unsigned int b = color & 0xFF;

    printf("R = 0x%02x\n", r);
    printf("G = 0x%02x\n", g);
    printf("B = 0x%02x\n", b);
}
\end{example}

\iosample{
}{
R = 0xc0 \\
G = 0x40 \\
B = 0x80 \\
}

\texttt{0xC04080} 的三個 byte 分別是 \texttt{0xC0}、\texttt{0x40}、
\texttt{0x80}。 \texttt{(color >> 16) \& 0xFF} 把 R 滑到最低 8 位再截
斷；右移 $8$ 位取出 G；不移位、直接 AND，取出本來就在最右邊的 B。 這
個位移加遮罩的組合，拆封包、拆機器字時會一再出現。

反過來更常寫：手上有三個 $0$ 到 $255$ 的分量，要組回一個整數。 每個
分量先左移到自己的區間，再用 OR 併起來。 AND 負責截斷，OR 負責拼接，
移位負責對準。

\begin{example}{c}{(pack-color.c) 組回 RGB}
#include <stdio.h>

int main(void) {
    unsigned int r = 0xC0;
    unsigned int g = 0x40;
    unsigned int b = 0x80;
    unsigned long color =
        ((unsigned long)r << 16) | ((unsigned long)g << 8) | b;

    printf("color = 0x%06lx\n", color);

    unsigned int r2 = (color >> 16) & 0xFF;
    unsigned int g2 = (color >> 8) & 0xFF;
    unsigned int b2 = color & 0xFF;
    printf("R = 0x%02x\n", r2);
    printf("G = 0x%02x\n", g2);
    printf("B = 0x%02x\n", b2);
}
\end{example}

\iosample{
}{
color = 0xc04080 \\
R = 0xc0 \\
G = 0x40 \\
B = 0x80 \\
}

左移之前先轉成 \texttt{unsigned long}，避免 $32$ 位 \texttt{int} 把高
位切掉。 組回去的值跟上一例拆開前相同，再拆一次仍得原來的三個分量。

\subsection{計數 1 的個數}

給定整數，數一數二進位表示中有多少個 1。

\begin{example}{c}{(count-bits.c) 計算 1 的個數}
#include <stdio.h>

int count_ones(unsigned int x) {
    int count = 0;
    while (x) {
        count += (x & 1);  // check lowest bit
        x >>= 1;           // right shift
    }
    return count;
}

int main(void) {
    printf("%d\n", count_ones(7));   // 111 = 3 ones
    printf("%d\n", count_ones(12));  // 1100 = 2 ones
}
\end{example}

\iosample{
}{
3 \\
2 \\
}

\subsection{進階：更快的計數技巧（選讀）}

上面的迴圈每一步只看最低位，再右移一格。 $1$ 若出現在很高的位，就得一
直移到那一格為止；最壞情況接近型別的位寬。 有一個常見的加速：
\texttt{n \&= (n - 1);} 會一次清掉最低那個 $1$，其餘位維持原樣。

\begin{code}{c}{一次清掉最低位的 1}
int count_ones_fast(unsigned int n) {
    int count = 0;
    while (n) {
        n &= (n - 1);   // clear the lowest 1-bit
        count++;
    }
    return count;
}
\end{code}

以 $12$（\texttt{1100}）為例。 $12 - 1 = 11$（\texttt{1011}），AND 之
後得 $8$（\texttt{1000}），最低的那個 $1$ 消失了。 下一步
$8 - 1 = 7$（\texttt{0111}），AND 得 $0$，迴圈結束。 總共走 $2$ 次，
剛好等於 $1$ 的個數。 前面那個右移版本對 $12$ 要走 $4$ 步才移完。

一般情況差在迴圈次數跟什麼成正比。 右移版每一步只看最低位，一定走到
最高那個 $1$ 為止，步數接近型別位寬；最壞是 $32$ 位裡只有 MSB 為 $1$
，仍要走滿 $32$ 步。 加速版每一步清掉一個 $1$，步數等於 $1$ 的個數。
\texttt{0x80000001} 只有兩個 $1$，一在 bit $0$、一在 bit $31$：右移版
走 $32$ 步，加速版走 $2$ 步。 高位含很多 $1$、中間又夾著長串 $0$ 時，
兩種寫法的差距就是位寬對上 popcount。

產品程式通常不必手寫這兩種迴圈。 GCC 與 Clang 提供
\texttt{\_\_builtin\_popcount}（64 位用 \texttt{\_\_builtin\_popcountll}
），編譯器會把它編成硬體的 popcnt 指令。 作業、講原理、或編譯器沒有
內建時，用上面的迴圈；其餘時候優先用內建。

\begin{code}{c}{編譯器內建的 popcount}
unsigned int x = 0x80000001u;
int n = __builtin_popcount(x);   // 2; GCC / Clang
\end{code}

\keypoint{移位運算在二進位層面上動作，所以超過型別位寬是未定義行為。
例如 32 位整數右移 32 格結果不可預測。}

\begin{drills}[label=drills:bit-shift]
\item \texttt{unsigned int x = 3;}，\texttt{x << 2} 的結果是多少？
\item 有號負數右移的結果，C 標準保證了什麼、又沒保證什麼？
\item 要從 \texttt{0xC04080} 取出中間那個位元組（G），該怎麼寫？
\item \texttt{n \&= (n - 1);} 這一行對 \texttt{n} 的位元模式做了什麼事？
\end{drills}

\section{綜合應用}
\label{sec:bit-applications}

切換第 $n$ 位就是 XOR 加上移位做出的遮罩，上一節的
\texttt{FLAG\_GREEN} 已經做過一次。 這裡用 Unix 風格的檔案權限，把五
個運算子放進同一支程式：OR 設置、AND 檢查、AND 加互補清除、XOR 切換、
移位負責把 owner／group／other 三組 $3$ 位對準。

每一組權限占 $3$ 位，意義與十進位 $4$、$2$、$1$ 相同：讀、寫、執行。
owner 放在 bit $8$ 到 $6$，group 在 bit $5$ 到 $3$，other 在 bit $2$
到 $0$。 \texttt{0755} 這種八進位字面值，剛好一位一位對準這 $9$ 格。

\begin{example}{c}{(file-mode.c) 檔案權限的位元}
#include <stdio.h>

#define R 4u
#define W 2u
#define X 1u

unsigned int pack(unsigned int owner, unsigned int group,
                  unsigned int other) {
    return (owner << 6) | (group << 3) | other;
}

int main(void) {
    unsigned int mode = pack(R | W | X, R | X, R | X);
    printf("mode = 0%o\n", mode);

    unsigned int owner = (mode >> 6) & 7u;
    printf("owner = %u\n", owner);

    if ((owner & W) != 0) {
        printf("owner can write\n");
    }

    mode &= ~(W << 6);
    printf("clear owner write: 0%o\n", mode);

    mode ^= X;
    printf("toggle other exec: 0%o\n", mode);
}
\end{example}

\iosample{
}{
mode = 0755 \\
owner = 7 \\
owner can write \\
clear owner write: 0555 \\
toggle other exec: 0554 \\
}

\texttt{pack} 裡的左移加 OR，跟組 RGB 是同一招，只是每段 $3$ 位而不是
$8$ 位。 \texttt{(mode >> 6) \& 7u} 取出 owner：右移對準，AND
\texttt{7u}（二進位 \texttt{111}）當擋板，只留 $3$ 位。
\texttt{(owner \& W) != 0} 是 AND 節的檢查。 \texttt{\~{}(W << 6)} 先
把寫權限對準 owner 那一段再翻反，AND 回去就清掉它，$0755$ 變成
$0555$。 最後 \texttt{mode \textasciicircum{}= X} 切換 other 的執行位
，變成 $0554$。

結構章會看到 \texttt{unsigned x : 3} 這種由編譯器切位的寫法。 本章先
練手動遮罩與移位；欄位較短、又對齊機器字時，手寫位元運算比較能掌握每
一格在哪。 欄位很多、只想依名字讀寫時，再讓編譯器切。

\subsection{優先級陷阱}

位元運算的優先級\textbf{低於}關係運算子。 括號很重要。

\begin{code}{c}{運算子優先序}
// wrong: a | b == c parsed as a | (b == c)
// correct
if ((a | b) == c) {
    printf("OK\n");
}

// wrong: x & 0x0F = 0 uses = not ==
// correct
if ((x & 0x0F) == 0) {
    printf("OK\n");
}
\end{code}

\keypoint{位元運算子順序：\texttt{\&} 高於 \texttt{\textasciicircum{}}
高於 \texttt{|}。 複雜表達式一定要括號。}

\subsection{慣用語總表}

這一章教過的核心寫法可以收成下面這張小抄。 各行都假設被操作的變數是
\texttt{unsigned}，遮罩寫成 \texttt{1u << n}。 從設置、清除到切換都守
這兩條，算術位移與符號位延伸就進不來。

\begin{code}{c}{位元慣用語小抄}
x |=  (1u << n);              // set bit n
x &= ~(1u << n);              // clear bit n
x ^=  (1u << n);              // toggle bit n
on  = (x & (1u << n)) != 0;   // test bit n
bit = (x >> n) & 1;           // read bit n
a ^= b; b ^= a; a ^= b;       // XOR swap (a and b distinct)
\end{code}

設置、清除、切換分別對應 OR、AND 加互補、XOR；測試與讀取對應 AND，差
在前者只問亮不亮、後者把那一位抽出來變成 $0$ 或 $1$。 交換那一行記得
避開同一個位址。 檔案權限那支程式把這張小抄全用過一遍：
\texttt{pack} 是移位加 OR，取出 owner 是移位加 AND，關寫權限是 AND 加
互補，切換執行位是 XOR。

\begin{drills}[label=drills:bit-applications]
\item Unix 檔案權限裡，owner 的讀寫執行三位放在哪個 bit 範圍？
\item \texttt{a | b == c}（沒加括號）會被剖析成什麼？ 該怎麼加括號修
正？
\item \texttt{mode \&= \~{}(W << 6);} 這一行在做什麼？
\end{drills}

\addcontentsline{toc}{section}{練習解答}
\section*{練習解答}

各節練習的答案收在這裡。 手算題直接給數值結果，觀念題給能一句話核對
的簡答。

\paragraph{第~\ref{sec:bit-basics}~節}
\begin{enumerate}
\item $13 = 8 + 4 + 1 = 2^3 + 2^2 + 2^0$，二進位是 \texttt{1101}。
\item MSB 是第 7 位（一個位元組的最高位），權重是 $2^7 = 128$。
\item \texttt{0x49} 是 \texttt{0100 1001}，十進位是 $73$。
\item 等價於 \texttt{x = x | (1u << 2)}。
\end{enumerate}

\paragraph{第~\ref{sec:bit-and}~節}
\begin{enumerate}
\item $147 \& 61 = 17$（\texttt{10010011 \& 00111101 = 00010001}）。
\item \texttt{(x \& (1u << 3)) != 0}，或等價地 \texttt{(x >> 3) \& 1}。
\item 因為 \texttt{flags} 上其他沒被 \texttt{MASK} 覆蓋到的位是 0 或
1 都不該影響判斷；先 AND 再比，只問 \texttt{MASK} 指定的那幾位。直接
比較兩個整數則要求每一位都相同。
\item 因為 \texttt{\&} 優先序低於 \texttt{==}，會先算
\texttt{MASK == MASK}（結果 1），再跟 \texttt{flags} 做 AND，跟原意完
全不同。
\end{enumerate}

\paragraph{第~\ref{sec:bit-or}~節}
\begin{enumerate}
\item $182 | 205 = 255$（\texttt{10110110 | 11001101 = 11111111}）。
\item \texttt{|=} 只把指定的位設成 1，其餘位不動；直接賦值
\texttt{flags = FLAG\_BLUE} 會把整個 \texttt{flags} 換成新值，原本亮
著的其他旗標全部被蓋掉熄滅。
\item 字面值 \texttt{1} 是有號 \texttt{int}，\texttt{1 << 31} 會把 1
推進符號位、行為未定義；\texttt{1u} 是 \texttt{unsigned}，位移結果穩
定明確。
\end{enumerate}

\paragraph{第~\ref{sec:bit-complement}~節}
\begin{enumerate}
\item $x = 8$（\texttt{1100 \& \~{}0100 = 1100 \& 1011 = 1000}）。
\item \texttt{\~{}0u} 的位寬會自動跟著變數的型別走；寫死
\texttt{0xFF}／\texttt{0xFFFF} 等於把遮罩寬度綁死在 8 位或 16 位，型
別更寬時高位會留下不該有的 0。
\item \texttt{flags \&= \~{}(FLAG\_GREEN | FLAG\_RED);}
\end{enumerate}

\paragraph{第~\ref{sec:bit-xor}~節}
\begin{enumerate}
\item $182 \oplus 205 = 123$（\texttt{10110110 \textasciicircum{}
11001101 = 01111011}）。
\item 因為同一位 XOR 兩次會抵銷：先翻一次變成不同，再翻一次翻回原
值，等於沒做。
\item 因為任何數 XOR 自己都是 0；若兩個運算元其實是同一格記憶體，第
一步就把資料清成 0，後面兩步只是 0 XOR 0，資料整個消失。
\item 第一次切換後 \texttt{FLAG\_GREEN} 的位被點亮，第二次切換同一位
再翻一次，回到原來的值——等於沒變。
\end{enumerate}

\paragraph{第~\ref{sec:bit-shift}~節}
\begin{enumerate}
\item $3 << 2 = 12$（\texttt{0011} 變 \texttt{1100}）。
\item C 標準只保證 \texttt{unsigned} 右移補 0（邏輯位移）；有號負數右
移是 implementation-defined，可能補 0 也可能複製符號位，不保證跨編譯
器一致。
\item \texttt{(color >> 8) \& 0xFF}。
\item 一次清掉 \texttt{n} 裡最低的那個 1 位元，其餘位維持不變。
\end{enumerate}

\paragraph{第~\ref{sec:bit-applications}~節}
\begin{enumerate}
\item bit 8 到 bit 6。
\item 會被剖析成 \texttt{a | (b == c)}，因為 \texttt{|} 優先序低於
\texttt{==}；要改成 \texttt{(a | b) == c}。
\item 把寫權限（\texttt{W}）對準 owner 那一段（左移 6 位）再翻反當擋
板，AND 回去清掉 owner 的寫入位，其餘位不動。
\end{enumerate}
